\chapter{Installation}\label{ch:install}

\NF{} is distributed as a Python package with a compiled C++ core. On the supported platforms
binary wheels are available, so no compiler is needed. This chapter describes installation from
PyPI, the optional components, building from source on the three supported operating systems,
running the test suites and setting up the local language model.

\section{Installing from PyPI}\label{sec:install-pypi}

From version 0.9.0, \NF{} is published on the Python Package Index:

\begin{lstlisting}[style=shell]
python -m pip install nurbforge              # kernel and Python API
python -m pip install "nurbforge[mesh]"      # plus gmsh for meshing and STEP via OCC
\end{lstlisting}

Use a virtual environment so that \NF{} and its dependencies do not interfere with other
projects:

\begin{lstlisting}[style=shell]
python3 -m venv .venv
source .venv/bin/activate          # Windows PowerShell: .venv\Scripts\Activate.ps1
python -m pip install --upgrade pip
python -m pip install "nurbforge[mesh]"
\end{lstlisting}

The only mandatory runtime dependency is NumPy ($\ge 1.24$). pip selects the matching wheel for
your interpreter and platform automatically:

\begin{tabularx}{\textwidth}{@{}lX@{}}
\toprule
\textbf{Wheel tag} & \textbf{Platform} \\
\midrule
\code{manylinux_2_28_x86_64} & Linux x86\_64 with glibc $\ge 2.28$: AlmaLinux, Rocky and RHEL 8, 9
  and 10, Debian $\ge$ 10, Ubuntu $\ge$ 18.10, Fedora and others. \\
\code{win_amd64} & 64-bit Windows 10 and 11. \\
\code{macosx_14_0_arm64} & macOS 14 (Sonoma) or newer on Apple Silicon. \\
\code{macosx_14_0_x86_64} & macOS 14 or newer on Intel Macs. \\
\bottomrule
\end{tabularx}

Wheels exist for CPython 3.10 to 3.14. On other platforms (for example Linux on ARM, or
musl-based distributions such as Alpine) pip falls back to the source distribution and compiles
the C++ core. You then need a C++20 compiler and CMake ($\ge 3.24$); see
\cref{sec:install-source}.

\subsection{Wheels from GitHub releases}
Every release on \url{https://github.com/PhotonPower/Nurbforge/releases} also carries all wheels
and the source distribution as assets. This is useful on machines without internet access:
download the wheel that matches the Python version (\code{cp312} = CPython 3.12) and platform,
copy it over and install the file directly:

\begin{lstlisting}[style=shell]
python -m pip install nurbforge-0.9.0-cp312-cp312-win_amd64.whl
\end{lstlisting}

NumPy (and gmsh, if wanted) must then be installed the same way or be already present.

\section{Optional components (extras)}\label{sec:install-extras}

\begin{tabularx}{\textwidth}{@{}lX@{}}
\toprule
\textbf{Extra} & \textbf{Installs and enables} \\
\midrule
\code{nurbforge[mesh]} & \code{gmsh>=4.11}: the \code{nurbforge.mesh} package (meshing, gmsh
  physical groups, BEM surface meshes, STEP/BREP/IGES export through OpenCASCADE). \\
\code{nurbforge[llm]} & Nothing extra: the Ollama client uses only the standard library. The
  extra exists so that installation instructions can name it. \\
\code{nurbforge[dev]} & pytest, gmsh, ruff and mypy: everything needed to run the test suite and
  the linters. \\
\code{nurbforge[docs]} & Sphinx, furo, MyST and Breathe for building the HTML documentation. \\
\bottomrule
\end{tabularx}

Several extras can be combined: \code{pip install "nurbforge[mesh,dev]"}. The quotes stop the
shell from interpreting the brackets.

\subsection{gmsh on Linux}\label{sec:install-gmsh-linux}
The gmsh wheel from PyPI is built with its graphical user interface. Even when used headless it
loads a few X11 and OpenGL libraries at import time. On minimal server installations and in
containers they are often missing. \code{import gmsh} then fails with an \code{OSError}, not an
\code{ImportError}, for example \code{libGLU.so.1: cannot open shared object file}. On AlmaLinux
and RHEL install:

\begin{lstlisting}[style=shell]
sudo dnf install -y mesa-libGL mesa-libGLU libglvnd-opengl libgomp fontconfig \
                    libX11 libXext libXfixes libXcursor libXinerama libXft libXrender
\end{lstlisting}

On Debian and Ubuntu the corresponding packages are \code{libglu1-mesa}, \code{libgl1},
\code{libxft2}, \code{libxcursor1}, \code{libxinerama1} and \code{libgomp1}. No display or X
server is needed. On Windows and macOS the gmsh wheel works as is.

\begin{tip}
Code that should run with and without gmsh catches both exceptions:
\code{except (ImportError, OSError)}. The examples in this manual and in the repository do so.
\end{tip}

\section{Checking the installation}\label{sec:install-check}

The following script prints the version, the compiled core module and the gmsh status, and
computes the volume of a unit sphere as a quick self-test:

\example{02_check_installation.py}

On macOS with Apple Silicon and Python 3.14 it prints, for example,

\begin{lstlisting}[style=text]
Nurbforge 0.9.0 on Python 3.14.5 (Darwin arm64)
compiled core: _core.cpython-314-darwin.so
unit sphere volume 4.188790204786
gmsh 4.15.2
\end{lstlisting}

The exit status is 0 if the volume agrees with $4\pi/3$ to $10^{-8}$. If \code{import nurbforge}
itself fails with \code{ImportError: ... _core}, the compiled module does not match the
interpreter. This typically happens after an editable install with one Python version and use with
another, or when the current directory is the repository root and the source tree shadows the
installed package. Reinstall, or run from another directory.

\section{Reference for coding agents}\label{sec:install-agent-reference}

The package contains a compact API reference written for language models and coding agents: a
single Markdown file (about 80\,kB, 20\,000 tokens) with the conventions (units, parameter
domains, orientation, labels, errors), a worked example and every public function, class,
property and LLM tool with signature and a one-paragraph description. An agent that can only
install the wheel, without access to this manual or the repository, finds its way with it:

\begin{lstlisting}[style=text]
python -m nurbforge > nurbforge-reference.md     # or: nurbforge.reference() in Python
\end{lstlisting}

The file (\code{REFERENCE.md} inside the package) is generated from the API like
\cref{app:reference} and is attached to every GitHub release as
\code{nurbforge-reference.md}.

\section{Building from source}\label{sec:install-source}

A source build is needed to work on \NF{} itself, to use the C++ library directly, or on platforms
without wheels. The repository is
\begin{lstlisting}[style=shell]
git clone https://github.com/PhotonPower/Nurbforge.git
cd Nurbforge
\end{lstlisting}

The build has two parts that can be used independently:
\begin{enumerate}
  \item the C++ core, unit tests and examples, built with CMake directly
    (\cref{sec:install-cmake});
  \item the Python package, built by \code{pip} through scikit-build-core, which runs the same
    CMake files and compiles the pybind11 module (\cref{sec:install-python-src}).
\end{enumerate}

\subsection{Requirements}

\begin{tabularx}{\textwidth}{@{}lX@{}}
\toprule
\textbf{Tool} & \textbf{Version and notes} \\
\midrule
C++ compiler & C++20: GCC $\ge$ 11 (AlmaLinux 10 ships GCC 14), MSVC from Visual Studio 2022
  or 2026, Apple Clang from the Xcode Command Line Tools (macOS $\ge$ 14). \\
CMake & $\ge 3.24$; on Windows with Visual Studio 2026 $\ge 4.2$. Installing it with
  \code{pip install cmake} into the virtual environment is the easiest way to get a recent one. \\
Ninja & Used by the \code{debug}, \code{release} and \code{ci} presets; \code{pip install ninja}. \\
Python & $\ge 3.10$ with development headers (Linux: \code{python3-devel}). \\
pybind11, scikit-build-core & Needed for the Python module; pip installs them automatically for
  \code{pip install .}. \\
Git & To clone the repository. Catch2 (C++ tests) is fetched by CMake if not installed. \\
\bottomrule
\end{tabularx}

\subsection{AlmaLinux 10 / RHEL 10}\label{sec:install-alma}
These steps mirror the \code{almalinux10} CI job, which runs them on every change in an
\code{almalinux:10} container.

\begin{lstlisting}[style=shell]
sudo dnf install -y gcc-c++ git python3 python3-devel python3-pip
# runtime libraries for the gmsh wheel (see above)
sudo dnf install -y mesa-libGL mesa-libGLU libglvnd-opengl libgomp fontconfig \
                    libX11 libXext libXfixes libXcursor libXinerama libXft libXrender
python3 -m venv .venv
source .venv/bin/activate
pip install --upgrade pip cmake ninja pybind11 scikit-build-core
\end{lstlisting}

CMake and Ninja come from pip, so the distribution versions do not matter.

\subsection{Windows 11}\label{sec:install-windows}
Install Visual Studio 2022 or 2026 with the workload \emph{Desktop development with C++} (the
Community edition is sufficient), Python $\ge 3.10$ (from python.org, \code{winget install
Python.Python.3.12} or uv) and Git. In PowerShell, from the repository root:

\begin{lstlisting}[style=shell]
python -m venv .venv
.venv\Scripts\Activate.ps1
python -m pip install --upgrade pip cmake ninja pybind11 scikit-build-core
\end{lstlisting}

The \code{msvc} preset uses the newest installed Visual Studio with its multi-configuration
generator. No \emph{Developer PowerShell} is needed, because CMake finds MSVC itself. The Ninja
presets (\code{debug}, \code{release}) need \code{cl.exe} on \code{PATH} and therefore have to run
from a \emph{Developer PowerShell for VS}.

\begin{warning}[Windows pitfalls]
\begin{itemize}
  \item \code{Python was not found}: the Microsoft Store alias is active but no Python is
    installed. Install Python or disable the alias under \emph{Settings $\to$ Apps $\to$ Advanced
    app settings $\to$ App execution aliases}.
  \item \code{Building for: NMake Makefiles} or \code{Running 'nmake' '-?' failed}: CMake is too
    old for the installed Visual Studio. Use the CMake from the virtual environment.
  \item After switching compilers or CMake versions, delete \code{build/msvc} and configure
    again.
\end{itemize}
\end{warning}

\subsection{macOS}\label{sec:install-macos}
macOS 14 or newer with the Xcode Command Line Tools (\code{xcode-select --install}) and Python
$\ge 3.10$ (python.org or Homebrew) is all that is needed:

\begin{lstlisting}[style=shell]
python3 -m venv .venv
source .venv/bin/activate
pip install --upgrade pip cmake ninja pybind11 scikit-build-core
\end{lstlisting}

\paragraph{Deployment target.} Without further settings, Apple Clang builds for the running
macOS version, and the binaries may not start on older systems. To build for other Macs, set
the minimum version before configuring:
\begin{lstlisting}[style=shell]
export MACOSX_DEPLOYMENT_TARGET=14.0
\end{lstlisting}
The release wheels use 14.0. The lower bound is 13.3, because from that version Apple's libc++
provides \code{std::to_chars} for \code{double}, which the STEP writer uses. Below 13.3 CMake
stops with an explicit message.

\subsection{Building the C++ core and tests}\label{sec:install-cmake}

The repository contains CMake presets. Build trees go to \code{build/<preset>}.

\begin{lstlisting}[style=shell]
cmake --preset release -DNF_BUILD_PYTHON=OFF      # Windows: --preset msvc
cmake --build --preset release -j
ctest --preset release
\end{lstlisting}

\begin{tabularx}{\textwidth}{@{}lX@{}}
\toprule
\textbf{Preset} & \textbf{Configuration} \\
\midrule
\code{debug} & Ninja, \code{CMAKE_BUILD_TYPE=Debug}. \\
\code{release} & Ninja, \code{CMAKE_BUILD_TYPE=Release}. \\
\code{ci} & Release with warnings as errors and the benchmark executable, as used in CI. \\
\code{msvc} & Windows only: newest Visual Studio, multi-configuration; build preset uses
  \code{Release}. \\
\bottomrule
\end{tabularx}

The CMake options are

\begin{tabularx}{\textwidth}{@{}llX@{}}
\toprule
\textbf{Option} & \textbf{Default} & \textbf{Meaning} \\
\midrule
\code{NF_BUILD_TESTS} & \code{ON} & Catch2 unit tests (\code{tests/cpp}). \\
\code{NF_BUILD_PYTHON} & \code{ON} & The pybind11 module (needs pybind11 and Python headers). \\
\code{NF_BUILD_EXAMPLES} & \code{ON} & The C++ examples in \code{examples/cpp}. \\
\code{NF_WARNINGS_AS_ERRORS} & \code{OFF} & Treat compiler warnings as errors. \\
\code{NF_BUILD_BENCHMARKS} & \code{OFF} & The benchmark executable (\code{benchmarks/}). \\
\bottomrule
\end{tabularx}

\code{cmake --install build/release --prefix <dir>} installs the static library, the headers
and a CMake package configuration, so other projects can use \code{find_package(Nurbforge)}
(\cref{ch:cpp}).

\subsection{Building the Python package}\label{sec:install-python-src}

\begin{lstlisting}[style=shell]
pip install -e ".[dev]"        # editable: Python sources are used from the checkout
pip install ".[mesh]"          # or a regular installation
\end{lstlisting}

scikit-build-core configures CMake with \code{SKBUILD} set. Tests, examples and the install rules
of the C++ library are then skipped, and only the core library and the module \code{_core} are
built. With an editable install, changes to the Python files take effect immediately. Changes
to C++ sources need another \code{pip install -e .}.

\section{Running the tests}\label{sec:install-tests}

\begin{lstlisting}[style=shell]
ctest --preset release                 # C++ unit tests (Catch2)
pytest                                 # Python tests (tests/python)
pytest tests/python/test_mesh.py -v    # a single module
ruff check python tests examples       # lint
\end{lstlisting}

The Python tests also run every example script in \code{examples/python} and
\code{docs/manual/examples}. Tests that need gmsh are skipped automatically if gmsh cannot be
imported or loaded. The LLM tests use a scripted fake client and need no running model.

\section{Setting up the language model (optional)}\label{sec:install-ollama}

The LLM interface (\cref{ch:llm}) talks to a model served locally by
\href{https://ollama.com}{Ollama}. Nothing else needs to be installed on the Python side.

\begin{enumerate}
  \item Install Ollama from \url{https://ollama.com/download} (Linux, Windows, macOS).
  \item Start the server and pull a model that supports tool calling:
\begin{lstlisting}[style=shell]
ollama serve &              # on Windows and macOS the app starts the server
ollama pull qwen2.5:14b     # the default model; llama3.1, mistral-nemo, ... also work
\end{lstlisting}
  \item Start the interactive shell or use the Python API:
\begin{lstlisting}[style=shell]
nurbforge-llm --model qwen2.5:14b
\end{lstlisting}
\end{enumerate}

The server address and the model are taken from the environment variables
\code{OLLAMA_HOST} (default \code{http://localhost:11434}) and \code{NURBFORGE_LLM_MODEL}
(default \code{qwen2.5:14b}); see \cref{app:cli}. A 14-billion-parameter model needs about 9\,GB
of memory. Smaller models work, but they make more mistakes in multi-step requests.
